\section{Why the constructed neural controls have invariant laws}
\label{app:bounded-network-proof}
The following argument concerns the explicitly forced, bounded-parameter
models in Appendix~\ref{app:constructed-vs-ordinary}; it does not assert
recurrence for ordinary-network SGD.

\begin{proposition}[Bounded data gradients and latent return]
\label{prop:bounded-network-recurrence}
On a finite-dimensional active parameter space, consider the exact update
\[
 \begin{gathered}
 \theta^+=a\theta+b_B(\theta)+\nu Z,\qquad
 a=1-\eta\mu,\quad |a|<1,\\
 \sup_{B,\theta}\|b_B(\theta)\|\le K<\infty,
 \qquad Z\sim N(0,I),\quad\nu>0.
 \end{gathered}
\]
Assume fresh iid batches and Gaussian vectors, independent of each other
and of the current state, and continuous $b_B$.
The chain has a unique attracting invariant probability, a smooth
strictly positive density, and every finite polynomial moment.
These conclusions apply to the constructed tanh/sigmoid networks,
including their frozen-output-bias ablation.
\end{proposition}

\begin{proof}
First verify the bounded-gradient hypothesis for the stated networks.
Let $d_0,\ldots,d_L$, with $d_L=1$, be fixed layer widths, and write
$W_\ell=w_\ell\tanh\Theta_\ell$, $b_\ell=\tanh v_\ell$,
$w_\ell=2/\sqrt{d_{\ell-1}}$.
Finite inputs and targets give bounds $X_*,Y_*$. Hidden activations
and their derivatives have absolute value at most one. Thus hidden
outputs satisfy $A_\ell\le1$, with $A_0=X_*$, and the linear output
satisfies $|f_\theta(x_i)|\le F=d_{L-1}w_L+1$.
Set $D_L=1$ and $D_\ell=d_{\ell+1}w_{\ell+1}D_{\ell+1}$.
Backpropagation and $|\tanh'|\le1$ give, for every latent coordinate,
\[
 |\partial_{\Theta_{\ell,jk}}f_\theta|
 \le D_\ell w_\ell A_{\ell-1},\qquad
 |\partial_{v_{\ell,j}}f_\theta|\le D_\ell.
\]
There are finitely many coordinates, so $\|\nabla f_\theta\|\le J<\infty$
uniformly. Hence every individual and averaged squared-loss data gradient
has norm at most $G=(F+Y_*)J$.
The sampled objective therefore gives $b_B=-\eta g_B$, $K=\eta G$,
and $\nu=\eta\sigma_{\rm eff}>0$ exactly.

Put $q=|a|$ and $V=1+\|\theta\|$. The triangle inequality gives
$PV\le qV+1-q+K+\nu\mathbb E\|Z\|$.
On $\|\theta\|\le R$, all conditional Gaussian means have norm at
most $qR+K$. Consequently their mixture density, on a target ball
$\|u\|\le r$, is bounded below by
\[
 (2\pi\nu^2)^{-P/2}
 \exp[-(r+qR+K)^2/(2\nu^2)]>0,
\]
where $P$ is the active dimension. This gives compact minorization,
Lebesgue irreducibility, and aperiodicity. Localizing the displayed
geometric drift to a sufficiently large ball, the Foster--Harris theorem
gives uniqueness and weighted total-variation convergence
\citep{hairermattingly2011}.
For each $p>0$, subadditivity ($p\le1$) or Young's inequality ($p>1$)
gives $P\|\theta\|^p\le q_p\|\theta\|^p+C_p$, with $q_p<1$.
Iteration from a finite start, convergence in law, and Portmanteau give
the invariant $p$th moment. Invariance expresses its density as a mixture
of translated Gaussians with covariance $\nu^2I$; uniformly bounded
Gaussian derivatives justify differentiation and imply smoothness and
positivity. Freezing a bias removes its coordinate: the same argument
holds on the remaining space, where the forcing is still full rank.
\end{proof}

For the reported $\mu=.05$ and $\eta\in\{.5,10\}$, the outer factors
are $.975$ and $.5$. The proof establishes recurrence independently of
any measured tangent-growth sign; it proves neither positive Lyapunov
growth nor a density mode count.
